% -*- mode: LaTeX; TeX-PDF-mode: t; -*-     # Configure emacs auctex viewer
\input{./econtexRoot.texinput}

\documentclass{\econtex}\newcommand{\texname}{Syllabus}
\usepackage{\econtexSetup}\usepackage{\econtexShortcuts}

\provideboolean{Post}     % Post handouts to web as part of compilation - SLOW
\provideboolean{Copy}     % Copy handouts to Handouts directory as part of compilation
\setboolean{Post}{false}  % Don't post unless explicitly run from shell with an option set to do so
%\setboolean{Post}{true}
\opt{postPublic}{\setboolean{Post}{true}}
\setboolean{Copy}{true}
\opt{Web}{\setboolean{Copy}{false}\setboolean{Post}{false}}

% Definitions unique to this paper (`\texname` defined above)
\usepackage{\texname} % processes \texname.sty

\begin{document}\large\bibliographystyle{\econtexBibStyle}
\begin{center}{{\tiny \url{https://ccarrollatjhuecon.github.io/Choice-Syllabus-Latest}}}\end{center}

\begin{verbatimwrite}{\texname.title}
Syllabus for Intertemporal Choice
\end{verbatimwrite}

\thispagestyle{empty}

\newenvironment{blockpar}{\par\begin{minipage}{\textwidth}
\setlength{\parskip}{.5\baselineskip plus 1pt minus 1pt}}{\end{minipage}}

\pagestyle{plain}
\thispagestyle{empty}

\begin{center}
{\tiny \href{https://ccarrollatjhuecon.github.io/Choice-Syllabus-Latest}{\texname} \hfill \today}

{\LARGE Macroeconomic Theory I: Intertemporal Choice}
\medskip\medskip\medskip\medskip

\smallskip

Professor \href{http://www.econ2.jhu.edu/people/ccarroll}{Christopher D. Carroll}

\medskip

\medskip\medskip
Office Hours: \href{https://calendar.google.com/calendar/selfsched?sstoken=UUNxMjFfNUtvakNSfGRlZmF1bHR8NjZhZDA0NmI3YmE5N2Y5N2Y3YjE2MzI5ZDg5YTYwYjQ}{Make An Appointment Yourself (read the ``rules'' link)}

\medskip\medskip\medskip
\small
Teaching Assistant:  \\
Yiran (Emma) Ma \texttt{<\href{mailto:yma116@jh.edu}{\texttt{yma116@jh.edu}}>}
\\ Required discussion section: Thursdays, 1:30--2:20 pm, Wyman W603.
\\ Office Hours: TBA
\normalsize
\medskip

\small
\href{https://app.paperpile.com/econark-related-tkUex~5rSTa26MXwWY5FtZQ/all}{Link}: Link to paperpile folder which contains all the readings for this course.

\end{center}

\begin{small}

If you're a Google Calendar user, you should be able to subscribe to the course calendar using the following Calendar ID (after logging into your Google account, copy and paste the text into the ``Other Calendars'' box in \href{http://calendar.google.com}{\texttt{http://calendar.google.com}}):

\medskip
\centerline{\texttt{jhuecon.org\_v2kso2hp1memk0vka1gomvorm8@group.calendar.google.com}}
\medskip

If you are not yet a Google Calendar user, you should become one; items that I post on the calendar will be viewed as having been distributed to the class.  If you insist on boycotting Google, however, you can see the calendar if you click \href{https://www.google.com/calendar/embed?src=jhuecon.org_v2kso2hp1memk0vka1gomvorm8\%40group.calendar.google.com\&ctz=America/New_York}{\bf HERE}

\medskip

\noindent\textbf{Where everything lives.} Anything posted on the course calendar counts as
distributed to the class; the calendar can also be followed from Outlook or Apple Calendar
via \href{https://calendar.google.com/calendar/ical/jhuecon.org_v2kso2hp1memk0vka1gomvorm8%40group.calendar.google.com/public/basic.ics}{this iCal address},
or simply \href{https://calendar.google.com/calendar/embed?src=jhuecon.org_v2kso2hp1memk0vka1gomvorm8%40group.calendar.google.com}{viewed}.

\begin{description}\setlength{\itemsep}{1pt}
\item[Lecture notes] the book \href{https://jhu-econ.github.io/intertemporal-choice}{\texttt{jhu-econ.github.io/intertemporal-choice}}; each section of this syllabus links to its part.
\item[This year's problem sets, exams, and class recordings] the private repository \href{https://github.com/ccarrollATjhuecon/Choice-2026}{\texttt{github.com/ccarrollATjhuecon/Choice-2026}}. Access is by invitation: send your GitHub username to the TA. Create the account with your JHU address and apply for \href{https://github.com/education/students}{GitHub Education} (free).
\item[Past exams, with answers] \href{https://github.com/ccarrollATjhuecon/Choice}{\texttt{github.com/ccarrollATjhuecon/Choice}}.
\item[Recordings and summaries] every class is recorded; recordings are in Panopto (sign in with your JHED; reach them through Canvas $\rightarrow$ \emph{Panopto Video}, or the links in the repository's \texttt{Classes/README.md}), and the Zoom AI summary and transcript of each class are in the repository's \texttt{Classes/} folder.
\item[Computation] the \href{https://econ-ark.org}{Econ-ARK} toolkit (\href{https://github.com/econ-ark/HARK}{HARK}) and its demonstration notebooks (\href{https://github.com/econ-ark/DemARK}{DemARK}), which also run in the browser with nothing installed at \href{https://econ-ark.org/materials}{\texttt{econ-ark.org/materials}}. Bring a laptop to every class. The TA runs a setup session in the first two weeks; installation instructions are distributed before it.
\item[Submitting problem sets] as stated in each problem set's README in the repository.
\item[Who to ask] anything about the running of the course --- repository access, submissions, the section, software, absences --- goes to the TA, who is copied on everything; the economics comes to me.
\item[Grading] final exam 50\%, midterm 25\%, problem sets 10\%, participation 15\%.
\end{description}

\medskip

\end{small}
\normalsize

\medskip\medskip

Theory and evidence about the dynamic behavior of households, firms, and the macroeconomy as a whole in the short and long run.  Begins with a thorough discussion of the consumption/saving problem of households, including the role of uncertainty, then moves to investment behavior of firms, including the relationship between financial markets and firm behavior.  General equilibrium models of firms and households combine to generate benchmark models of economic growth, which leads us to a benchmark specification for dynamic aggregate macroeconomic models.

This syllabus contains required and recommended readings for each
topic.  Required readings are indicated by a ${*}$. 

You should own a copy of the texts by \cite{blanchard&fischer:text}, \cite{deatonUnderstandingC}, 
\cite{romer:text} and \cite{lsRecurse}.

{\small Students from other departments are welcome to take the course, but 
they should be prepared for regular use of higher mathematics (multivariate
calculus, probability and statistics, functional analysis) which will be employed
later in the course.}

\medskip\medskip
\noindent Course handouts are available at: 
\medskip

\texttt{\href{https://jhu-econ.github.io/intertemporal-choice}{jhu-econ.github.io/intertemporal-choice}}

\pagebreak


\begin{verbatimwrite}{./Sections/Method/title.tex}
Methodology for Macroeconomists
\end{verbatimwrite}
\section{Course Intro: Methodology for Macroeconomists}

\begin{verbatimwrite}{./Sections/Method/body.tex}

\begin{itemize}
\item[Bib:] \href{http://www.econ2.jhu.edu/people/ccarroll/courses/choice/Syllabus/Method.bib}{Method.bib}
\item[Handouts:]  \texttt{\href{https://www.econ2.jhu.edu/people/ccarroll/Public/Lecturenotes/CourseIntro/Methodology}{Methodology}}
\item[Readings:]
\end{itemize}

% \newcommand{\si}{\cite{summersIllusion}}
% \newcommand{\ac}{\cite{acemogluCrisis}}
% \newcommand{\gk}{\cite{grantKeynes}}
% %\newcommand{\kt}{\cite{krugmanThoughts}}
% \newcommand{\kh}{\cite{krugmanHistory}}
% \newcommand{\cp}{\cite{caballeroPretense}}
% \newcommand{\rr}{\cite{rodrikRules}}
% \newcommand{\bf}{\cite{barberaFight}}     

\begin{itemize} 
\reqd \cite{summersIllusion}  
\reqd \cite{acemogluCrisis}   
%\reqd \cite{keynesMaster}      
%\reqd \cite{krugmanThoughts} 
\reqd \cite{krugmanHistory}   
\reqd \cite{caballeroPretense}
%\reqd \cite{rodrikRules}      
\reqd \cite{barberaFight}     
\reqd \cite{romerTrouble}
%\reqd \cite{kocherlakotaTrouble}
\end{itemize}

\end{verbatimwrite}

\input ./Sections/Method/body.tex

\begin{verbatimwrite}{./Sections/Consumption/title.tex}
Consumption
\end{verbatimwrite}

\section{Consumption}\label{sec:consumption}
\begin{verbatimwrite}{./Sections/Consumption/body.tex}

\begin{itemize}
\item[Bib:] \texttt{\href{http://www.econ2.jhu.edu/people/ccarroll/courses/Choice/Syllabus/Consumption.bib}{Consumption.bib}}
\item[Handouts:] \texttt{\href{https://jhu-econ.github.io/intertemporal-choice/content/consumption}{Consumption}}
\item[Readings:]
\end{itemize}

\newcommand{\IntroC}{Introduction and Overview}
\subsection{\IntroC}\label{subsec:\IntroC}

\newcommand{\dt}{\cite{deatonUnderstandingC}}
\newcommand{\fr}{\cite{friedmanATheory}}
\bi
\reqd \fr, Chapters I, II, III
\reqd \dt, Preface; pp. 76-81
\ei

\subsection{The Life Cycle and Overlapping Generations Models}

\providecommand{\blf}{\cite{blanchard&fischer:text}}
\providecommand{\ro}{\cite{romer:text}}
\bi
\reqd \dt, pp. 1-6
\reqd \blf, Chapter 3: Introduction, pp. 91-113
\recm \cite{diamond:olg}
\recm \cite{feldstein:induced}
\reqd \cite{modigliani:nobel}
\recm \cite{summersCapTax}
\recm \cite{pemberton:failure}
\reqd \cite{carroll&summers:cparallelsy}
\ei

\subsection{The Infinite Horizon Representative Agent Model}

\bi
\reqd \dt, Chapter 3
\reqd \cite{hallRandomWalk}
\recm \cite{flavinSensitive}
\recm \cite{cdSmooth}
\reqd \cite{cmModel}
\recm \cite{hallSubstitution}
\recm \cite{cfwSentiment}
\ei

\subsection{Consumption with Uncertainty and/or Liquidity Constraints}

\bi
\reqd \dt, Chapter 6
\recm \cite{zeldes_jpe89}
\recm \cite{caballero:jme}
\recm \cite{deatonLiqConstr}
\recm \cite{toche:urisk}
\reqd \cite{carroll:atheorynberwp}
\reqd \cite{CarrollKimballPSPW}
\reqd \cite{carroll&summers:cparallelsy}
\recm \cite{BufferStockTheory}
\recm \cite{carroll:death}
\ei

\subsection{Consumption Based Asset Pricing}

\bi
\reqd \cite{lucas:assetpricing}
\reqd \cite{mehraPrescottPuzzle}
\recm \cite{blanchard:burstingbubbles}
\recm \cite{carrollIrrational}
\ei
\end{verbatimwrite}

\input ./Sections/Consumption/body.tex


\begin{verbatimwrite}{./Sections/Investment/title.tex}
Investment
\end{verbatimwrite}

\section{Investment}\label{sec:Investment}
\begin{verbatimwrite}{./Sections/Investment/body.tex}
\begin{itemize}
\item[Bib:] \texttt{\href{http://www.econ2.jhu.edu/people/ccarroll/courses/Choice/Syllabus/Investment.bib}{Investment.bib}}
\item[Handouts:]  \texttt{\href{https://jhu-econ.github.io/intertemporal-choice/content/investment}{Investment}}
\item[Readings:]
\end{itemize}

\providecommand{\blf}{\cite{blanchard&fischer:text}}
\providecommand{\ro}{\cite{romer:text}}

\bi
\reqd \blf, Chapter 2, pages 58--69; Chapter 6, pp 291--302
\recm \cite{hall&jorgenson:i}
\reqd \cite{hayashi:q}
\reqd \cite{abel&eberly:unified}
\reqd \cite{chh:reconsider}
\ei

\end{verbatimwrite}

\input ./Sections/Investment/body.tex

\begin{verbatimwrite}{./Sections/Growth/title.tex}
Growth
\end{verbatimwrite}

\section{Growth}\label{sec:Growth}

\begin{verbatimwrite}{./Sections/Growth/body.tex}

\providecommand{\blf}{\cite{blanchard&fischer:text}}
\providecommand{\ro}{\cite{romer:text}}
\begin{itemize}
\item[Bib:] \texttt{\href{http://www.econ2.jhu.edu/people/ccarroll/courses/Choice/Syllabus/Growth.bib}{Growth.bib}}
\item[Handouts:]  \texttt{\href{https://jhu-econ.github.io/intertemporal-choice/content/growth}{Growth}}
\end{itemize}
\bi
\reqd \blf, Chapter 2, pp 37-57
\reqd \cite{deatonStateCapacity}
\reqd \cite{salaimartin:lecnotes}
\recm \cite{king&rebelo:trans}
\recm \cite{phelps:golden}
\ei

\end{verbatimwrite}

\input ./Sections/Growth/body.tex

\begin{verbatimwrite}{./Sections/DSGEModels/title.tex}
Real Business Cycle Models
\end{verbatimwrite}

\section{Real Business Cycle Models}
\begin{verbatimwrite}{./Sections/DSGEModels/body.tex}
\begin{itemize}
\item[Bib:] \texttt{\href{http://www.econ2.jhu.edu/people/ccarroll/courses/Choice/Syllabus/DSGEModels.bib}{DSGEModels.bib}}
\item[Handouts:]  \texttt{\href{https://jhu-econ.github.io/intertemporal-choice/content/dsge}{DSGEModels}}
\end{itemize}

\providecommand{\blf}{\cite{blanchard&fischer:text}}
\providecommand{\ro}{\cite{romer:text}}
\bi
\reqd \blf, Chapter 1
\reqd \blf, Chapter 7
\reqd \cite{prescottTheoryAhead}
\reqd \cite{summers:skeptical}
\ei

\end{verbatimwrite}

\input ./Sections/DSGEModels/body.tex

\pagebreak

\href{https://app.paperpile.com/econark-related-tkUex~5rSTa26MXwWY5FtZQ/all}{Link}: Link to paperpile folder which contains all the readings for this course.

\renewcommand\refname{\noindent{\href{http://www.econ2.jhu.edu/people/ccarroll/courses/Choice/Syllabus.bib}{\large Readings {\small (click to download \texttt{.bib} file)}}}}

% Make the bibliography
% \input{\LaTeXInputs/bibliography_blend}
\bibliography{system}

\end{document}

